\documentclass{article}
\usepackage{amsmath,amssymb}
\usepackage[margin=1in]{geometry}
\usepackage{xurl}
\usepackage[hidelinks]{hyperref}
% Shared commands for notes in docs/paper-gaps/.

\DeclareUnicodeCharacter{2080}{\ifmmode{}_{0}\else\textsubscript{0}\fi}
\DeclareUnicodeCharacter{2081}{\ifmmode{}_{1}\else\textsubscript{1}\fi}
\DeclareUnicodeCharacter{2082}{\ifmmode{}_{2}\else\textsubscript{2}\fi}
\DeclareUnicodeCharacter{2099}{\ifmmode{}_{n}\else\textsubscript{n}\fi}

\newcommand{\Lean}{\textsc{Lean}}
\ExplSyntaxOn
\NewDocumentCommand{\leanid}{m}{
  \ifmmode
    \textsf{\detokenize{#1}}
  \else
    \group_begin:
    \urlstyle{sf}
    \tl_set:Nn \l_tmpa_tl {#1}
    \tl_replace_all:Nnn \l_tmpa_tl {\_} {_}
    \exp_args:NV \path \l_tmpa_tl
    \group_end:
  \fi
}
\ExplSyntaxOff
\providecommand{\tr}{\operatorname{tr}}
\newcommand{\tnleanIssueBase}{https://github.com/LionSR/TNLean/issues/}
\newcommand{\tnleanPrBase}{https://github.com/LionSR/TNLean/pull/}
\newcommand{\tnleanDocBase}{https://sirui-lu.com/TNLean/docs/}
\newcommand{\ghissue}[1]{\href{\tnleanIssueBase#1}{GitHub issue \##1}}
\newcommand{\ghpr}[1]{\href{\tnleanPrBase#1}{PR \##1}}
\newcommand{\doclink}[1]{\href{\tnleanDocBase#1.html}{\path{#1}}}

% Dirac notation and the identity operator, as in blueprint/src/macros/common.tex.
% \providecommand keeps note-local definitions authoritative.
\usepackage{dsfont}
\providecommand{\ket}[1]{|#1\rangle}
\providecommand{\bra}[1]{\langle #1|}
\providecommand{\braket}[2]{\langle #1 | #2 \rangle}
\providecommand{\ketbra}[2]{|#1\rangle\!\langle #2|}
\providecommand{\Id}{\mathds{1}}
\providecommand{\one}{\mathds{1}}
\providecommand{\C}{\mathbb{C}}
\providecommand{\MN}[1]{M_{#1}(\C)}
\providecommand{\MD}{M_D(\C)}

% Verdict marker: \gapnote{<kind>}{<status>}. Kind is the policy
% classification (clarification, local-correction, scope-restriction,
% unfaithful, false-source, open-gap); status is open, wip (elimination
% underway), resolved, or historical. Machine-read by the site index;
% prints a verdict line.
\newcommand{\gapnote}[2]{\par\noindent\textbf{Verdict:} #1 (#2).\par}

\title{Uniform native-closure regular PEPS renormalization}
\author{}
\date{2026-10-05}
\begin{document}
\maketitle
\gapnote{scope-restriction}{open}

\section{Source and proved specialization}
The source is Schuch, Cirac, and P\'erez-Garc\'ia,
\href{https://arxiv.org/abs/1001.3807v3}{arXiv:1001.3807v3},
Definition 5.6, Theorem 5.9 and Observation 6.6, local source lines
1515--1525, 1582--1621 and 1888--1909.

Let $G$ be an arbitrary finite group and $A$ a homogeneous regular
G-isometric four-leg tensor with an arbitrary finite physical alphabet.
For coarse periods $w,h\geq3$, one isometry $I$ on the full product block
physical support satisfies, for every native closure pair $(g_0,h_0)$,
\[
 I\Psi_f(g_0,h_0)=R\Psi_c(g_0,h_0)\otimes\Omega,\qquad
 I\widehat\Psi_f(g_0,h_0)=\widehat\Psi_c(g_0,h_0)\otimes\Omega.
\]
Here the fine periods are $2w,2h$, the coarse tensor is literally $A$,
$\Omega$ is a fixed unit Bell product, and
\[
 R=\left(\prod_v\frac{\sqrt{c_{B,v}}}{\sqrt{c_{C,v}}}\right)
    (\sqrt{|G|})^{|E(\Gamma)|}>0.
\]
Neither $I$, its local physical matrices, nor $R$ depends on the closure.
The factors are derived from the G-isometry of the actual four-site block
and of the original tensor with surplus identity registers. The previous
four-site Gram result retains its internal-cycle group-order factor.
Each native closure is proved nonzero rather than assumed nonzero.

The horizontal seam carries $h_0$ on the receiving left legs; the
vertical seam carries $g_0$ in the native downward orientation. Reversing
an arrow to express it as an ordered graph edge transposes its inserted
matrix and therefore inverts the regular group label. Both paired bond
labels receive the same insertion before the relative-coordinate change.
The identity $(ua)^{-1}(ub)=a^{-1}b$ leaves the surplus label unchanged.

\section{Coherent vectors and ground-state scope}
For arbitrary complex coefficients $z_{g_0,h_0}$, the same map gives
\[
 I\sum_{g_0,h_0}z_{g_0,h_0}\Psi_f(g_0,h_0)
 =R\left(\sum_{g_0,h_0}z_{g_0,h_0}\Psi_c(g_0,h_0)\right)\otimes\Omega.
\]
Fine and coarse sums vanish simultaneously, and their nonzero normalized
vectors satisfy the corresponding exact identity. Thus relative
coefficients and phases are preserved. Restricting the coefficients to
commuting pairs proves this assertion for arbitrary coherent vectors in
the commuting-closure span. The algebraic result also covers noncommuting
pairs but makes no parent-ground-space claim for them. Identification with
a particular parent kernel uses the separate parent-ground-space theorem;
it is not inferred from noncommuting closure nonvanishing.

\section{Remaining restrictions}
This result removes the identity-closure restriction of the earlier
periodic original-tensor fixed-point specialization. It preserves its
homogeneous, finite regular, physical-support scope. The native geometric
regrouping includes all positive coarse periods, but the physical-support
and Bell-separation argument requires the simple coarse graph and therefore
periods at least three. Extending it to bond-indexed loops and parallel
bonds is still required for periods one and two.

All original physical registers are retained. An ambient-space unitary
extension is not asserted; unused physical directions are permitted.
Open or other boundary conditions, nonregular representations, approximate
schemes, and a general parent-Hamiltonian renormalization identity remain
outside the theorem. These are the remaining qualifications to the source's
unrestricted fixed-point formulation.

\section{Declarations}
The scalar-one native geometric identity is
\leanid{TNLean.PEPS.torusGClosure_leftRegular_eq_twoByTwoBlocked}.
The native-to-bundled contraction comparison is
\leanid{TNLean.PEPS.torusGClosure_leftRegular_eq_twoByTwoBundledGraphSite}.
Uniform inserted physical support transport is
\leanid{TNLean.PEPS.exists_graphGIsometricInsertedPhysicalSupportTransport}.
The exact original-tensor surplus factorization is
\leanid{TNLean.PEPS.graphInsertedBondNetwork_regularGraphSurplusSite}.
The uniform capstone, including all coherent sums, is
\leanid{TNLean.PEPS.exists_regularTwoByTwoTwistedTensorIsometry}.
\end{document}
